% --- Patron: Ibeji (The Twin Stones) ---
\subsection{Ibeji — The Twin Stones (Partnership, Balance, Union)}

\textit{Lore.} Ibeji are twin spirits — one of stone, one of gear — who dance together to keep the walking cities of Ngomebe alive. They were not born; they were carved. The first stone city needed a keystone and a bearing, and the first Mason‑Kin and the first Aeler vent‑prior worked together for nine years to make them fit. When the city walked, the two spirits emerged from the joint — one from the stone, one from the gear — still holding hands. They teach that no craft is complete without another, that a bridge needs two shores, and that the strongest keystone is the one that bears the weight of a partner’s trust.

Among the peoples of Fate's Edge, Ibeji are invoked by masons, engineers, and any two people who must work in perfect trust. In the Mountain Satrapies, they are carved on the lintels of Aeler‑Dhaharan counting‑houses. In Ngomebe, they are painted on the rollers of every new city. Wherever a stone is set by one hand and sealed by another, Ibeji dance.

\begin{quote}
``Ibeji taught us that a single stone is just a rock. Two stones, set together, can hold back a river. One gear is a toy. Two gears can move a city. Find your twin. Work. Do not let go.''
\end{quote}

\noindent\textbf{Domain Focus}
\begin{itemize}
  \item \textbf{Partnership:} two hands, two crafts, two wills working as one
  \item \textbf{Balance:} equal weight, equal credit, equal sacrifice
  \item \textbf{The Union of Aeler and Ngomebe:} stone and gear, mountain and plateau
  \item \textbf{The Twin Keystone:} the joint that bears the load; break one, and both fall
\end{itemize}

\subsubsection*{Thiasos and Codex (Runekeeper Options)}
A Runekeeper sworn to Ibeji does not bind a single familiar — they bind \textbf{two}, for Ibeji are never one. Their \textbf{Thiasos} is a pair of small creatures that embody partnership: two twin mice that sleep curled together, two small lizards that change colour in synchrony, two beetles that roll the same ball of dung, two grey doves that nest in the same bell tower, or two small brass gears that click in perfect rhythm (these are not alive, but they hum when the partner is near). Each creature requires no food in the usual sense, but they must be kept together; separating them for more than a day makes them \textit{Compromised} (they grow listless, and the bearer suffers –1 die to all rolls involving coordination or assistance until reunited).

The \textbf{Codex} of Ibeji is never a book written by one hand. It may be a pair of matched slate tablets, each inscribed with half of every Rite; a single scroll that must be held by two people to be read; a set of twin bells that ring only when struck together; or a keystone with two faces, each carved by a different mason. Because Ibeji’s truths are written in duality, the Codex requires \textit{upkeep}: the Runekeeper must spend an hour each week with a partner (any willing person), reading each Rite aloud together, or the words become mismatched (treat as Compromised until re‑paired).

\subsubsection*{Patron's Gift (Imbuement)}
Once per scene as an action, a Runekeeper of Ibeji may touch two held items (weapons, tools, or pieces of armour) and whisper a word of union. For the rest of the scene, each item grants \textbf{+1 die} to rolls made to assist the other bearer, and the first time each scene the two bearers would be separated by force (a closing door, a collapsing floor, a spell of banishment), they may reroll the resistance check (the twin stones hold). After the scene, each user marks \textbf{+1 Obligation} as the joint tightens.

\subsubsection*{Rites of Ibeji}

\paragraph{Rite of the Balanced Weight (Low, 4 XP)} \hfill \\
\emph{Scene; Self and One Ally; No} \\
\textbf{Tags:} \texttt{[ASSIST]}, \texttt{[BOOST]}, \texttt{[BOND]} \\
\textbf{Materials:} Two stones of equal weight, or a single gear that has been split and re‑joined. \\
\textbf{Effect:} You and an ally within Near gain \textbf{+1 die} to any roll made to assist each other for the scene. When one of you takes the Assist action, the other gains an additional +1 die (stacking with the normal assist bonus). \\
\textbf{Push It:} The bond is visible — a faint line of golden light connects you. Any enemy who sees it must test \textit{Resolve} (DV 3) or be unable to target only one of you; they must attack both or neither. Mark \textbf{1 SB (Diamonds)} as the light draws attention. \\
\emph{Requires: Familiar (\textit{Invoke:} 1 Boon).} \\
\textbf{Invoke:} 1 action; mark \textbf{+1 Obligation}.

\paragraph{Rite of Ibeji’s Seal (Standard, 8 XP)} \hfill \\
\emph{Instant; Touch; No} \\
\textbf{Tags:} \texttt{[BOND]}, \texttt{[FATE]}, \texttt{[PROTECT]} \\
\textbf{Materials:} Two objects that belong to two different people (a ring, a lock of hair, a personal seal). \\
\textbf{Effect:} Touch two objects; for one day, they cannot be separated by less than magical force. If one is moved, the other moves with it. This can be used to lock a door (the door and its frame), to bind a prisoner’s hands (left and right manacle), or to keep two allies from being pulled apart by a trap. \\
\textbf{Push It:} The seal is permanent until broken by a ritual of separation (requiring the consent of both parties). Mark \textbf{+2 Obligation}. If one party breaks the seal by force, they suffer 2 Harm (spiritual backlash). \\
\emph{Requires: Familiar + Codex (\textit{Invoke:} 1 Boon).} \\
\textbf{Invoke:} 1 action; mark \textbf{+1 Obligation}.

\paragraph{Rite of the Dance of Stone (High, 14 XP)} \hfill \\
\emph{Extended; Zone; No} \\
\textbf{Tags:} \texttt{[MOVEMENT]}, \texttt{[CRAFT]}, \texttt{[SACRIFICE]}, \texttt{[DEBT]} \\
\textbf{Materials:} A keystone chip from a walking city, and a gear from an Aeler bearing. \\
\textbf{Effect:} You and an ally may move a stone city’s roller 1 mile without oxen — but you both mark 2 Fatigue. The city will remember your aid; the Mason‑Kin will grant you a \textit{Keystone Charter} (right to dwell in any Ngomebe city for a season) if you succeed. If you fail, the roller cracks, and the city stops until a Mason‑Kin can repair it. \\
\textbf{Push It:} Move the roller 2 miles instead of 1. Mark \textbf{+3 Obligation} and both you and your ally gain a permanent scar — a crack on your palm that never heals, matching the crack in the roller. \\
\emph{Requires: Familiar + Codex + Tier III (\textit{Invoke:} 2 Boons).} \\
\textbf{Invoke:} Extended rite; mark \textbf{+2 Obligation}. \\
\emph{Obligation:} 8 segments.

\subsubsection*{Cantors \& Cults of Ibeji}

\textbf{The Cantor’s Song.} A Cantor of Ibeji sings in \textit{duets} — melodies that require two voices to complete. One voice begins a phrase; the other finishes it. A Cantor alone can only hum half a tune. Their voice is used to synchronise work crews, to heal the rift between feuding partners, or to remind two people that they were once one team. Ibeji’s Cantors are often found on construction sites, in counting‑houses where two clerks share a ledger, or at weddings where they sing the vows as a round.

\textbf{The Cult of the Twin Keystone.} Ibeji’s cults meet in pairs — literally. Each member must bring a partner. The \textit{Twin Keystone} is a fellowship of masons, engineers, and Aeler‑Dhaharan families who have intermarried. A ritual involves each pair carving a single rune on a shared stone — one carves the left half, the other the right. If the rune is complete, the stone is added to a growing cairn. If it is mismatched, the pair must spend a day working together before they may try again. Their creed is carved on a single stone that is never split: \textit{“One stone is a rock. Two stones are a foundation. A foundation holds the world.”} Outsiders are welcome if they bring a partner. If they come alone, they will be given a partner — and the partner may not be human.

\subsubsection*{Witchcraft of Ibeji (Craft / Partnership)}

A witch who walks with Ibeji is a \textit{joint‑smith} or \textit{bond‑weaver} — one who crafts the places where two things meet. Their magic works through the tools of union: a double‑headed hammer that strikes two nails at once, a pair of tongs that cannot be separated by less than a forge’s heat, or a spirit level that has two bubbles — one for each partner. Their tools are never single; a hammer without a matching pair is just a club.

\textbf{The Witch’s Tool: The Twin Needles.} Two needles of cold iron, bound by a single thread that never tangles. Once per session, the witch may sew two objects together (two pages of a contract, two door handles, two sleeves of a shirt). For the next day, they cannot be separated without the witch’s consent. Cutting the thread causes both objects to rust instantly. The needles can also be used to stitch a wound closed; the patient and the healer share the pain (each takes 1 Fatigue, but the wound heals 1 Harm).

\textbf{A Hedge Gift: Shared Breath (2 XP).} Once per scene, you and an ally within Near may each spend 1 Boon to share a single breath. For the rest of the scene, whenever one of you rolls a 1, the other may take that Story Beat instead of the GM (the GM does not gain it). The beat must be spent on a complication that affects the character who took it.

\subsubsection*{Ibeji’s Corruption Manifestations}

\begin{tblr}{
  colspec = {Q[c,1.8cm] X[2.5] X[2.5]},
  rowhead = 1,
  row{odd} = {bg=gray!10},
  row{1} = {bg=gray!30, font=\bfseries},
  hlines,
}
Level & Benefit & Cost / Quirk \\
1 & \textbf{Steady Hand:} +1 die to Craft rolls when working with a partner. & \textbf{Partner‑Dependent:} –1 die to any Craft roll made alone. \\
2 & \textbf{Twin’s Instinct:} Once per scene, you may know the emotional state of your bonded partner without speaking. & \textbf{Shared Pain:} When your partner takes Harm, you must test \textit{Resolve} (DV 3) or mark 1 Fatigue. \\
3 & \textbf{Two as One:} +2 dice to assists, and the first assist each scene does not cost a Boon. & \textbf{Jealous Joint:} –1 die to social rolls with anyone your partner distrusts. \\
4 & \textbf{Unbreakable Seal:} Once per session, declare a bond between two objects or people; it cannot be broken by non‑magical means for one week. & \textbf{Compulsive Union:} You must try to reconcile any two arguing parties you witness; ignoring a quarrel costs 1 Fatigue. \\
5 & \textbf{Dance of Stone:} Once per session, you and a partner may move a massive object (a collapsed beam, a fallen statue) as if it weighed nothing — but you both mark 1 Fatigue. & \textbf{Cracked Keystone:} Each use of this power causes a visible crack to appear on your skin (like dried earth). After nine cracks, you become a living keystone — you cannot be separated from your partner by any means, but you also cannot leave their side. \\
6+ & \textbf{Avatar of the Twin Stones:} Once per session, you and a bonded partner may act on the same turn, sharing a single pool of actions (one Action each, but you may move as one). & \textbf{Twin’s Doom:} Mark +3 Obligation. If one of you dies, the other must test \textit{Resolve} (DV 6) or immediately fall into a catatonic state for 1d4 days. Your reflections now always appear in pairs, even when you are alone. \\
\end{tblr}

\subsection*{Playstyle Notes}
Ibeji rewards those who work with a partner, who share credit, and who understand that the strongest structure is the one that distributes the load. Followers become masters of cooperation, binding, and the art of making two things work as one. The corruption progression leads toward a being who can no longer function alone — but who, with the right partner, can move mountains. Ideal for married couples in character, for engineers and masons, and for any two players who want their characters to be more effective together than apart.

\subsection*{Rivalries \& Obligations}

\textbf{Major Rivalries:}
\begin{itemize}
  \item \textbf{The Clockwork Monad (Iteration):} The Monad perfects the individual; Ibeji perfects the pair. Their followers disagree on whether progress comes from better parts or better joints.
  \item \textbf{Solitary patrons (e.g., the Lone Hunter archetypes):} Ibeji’s followers distrust those who refuse to work with others.
\end{itemize}

\textbf{Hard Obligation Triggers:}
\begin{itemize}
  \item \textbf{7+ Segments:} You must settle a dispute between two craftsmen who have stopped speaking. The Hollow will witness the reconciliation; if it fails, both craftsmen will blame you.
  \item \textbf{10+ Segments:} Ibeji demands you choose a permanent partner — a bond that cannot be broken by any means. If you refuse, you lose the ability to receive assistance from anyone; all assistance rolls against you are Desperate. If you accept, you gain the \textit{Twin’s Doom} condition (see table) but also the certainty that you will never work alone again.
\end{itemize}

\noindent\textbf{Emphasises}
\begin{itemize}
  \item \textbf{Partnership:} two hands, two crafts, two wills as one
  \item \textbf{Balance:} equal weight, equal credit, equal sacrifice
  \item \textbf{Union of Aeler and Ngomebe:} stone and gear, mountain and plateau
  \item \textbf{The Twin Keystone:} the joint that bears the load; break one, and both fall
\end{itemize}